\subsection{局部紧域的结构}

域 Q 关于 Q 上的非真绝对值的完备化 R 和 \(Q_{p}\) 是局部紧的（p 为任意素数）。另一方面，对于素数 p 的每个幂 \(q = p^{f}\)，形式幂级数域 \(\mathbf{F}_{q}((\mathrm{T}))\)（在有限域 \(F_{q}\) 上）的拓扑由 §3 第 4 节例 3 中定义的赋值给出，因而是局部紧的：因为赋值环 \(F_{q}[[T]]\) 的极大理想由 T 生成；我们知道这个环关于 (T)-进拓扑完备（第三章，第 2 节，第 6 节，命题 6），又由于剩余域 \(F_{q}\) 是有限的，§ 5 第 1 节的命题 2 证明了我们的断言。反之：

定理 1. 设 K 为一个（不必交换的）非离散局部紧域。

\begin{enumerate}
\def\labelenumi{(\roman{enumi})}
\item
  若 \(\mathbf{K}\) 的特征为 0 且 \(\mathrm{mod}_{\mathbf{K}}\) 不是超度量绝对值，则 \(\mathbf{K}\) 同构于域 \(\mathbf{R}\)、\(\mathbf{C}\) 或 \(\mathbf{H}\) 之一。
\item
  若 \(\mathbf{K}\) 的特征为 0 且 \(\mathrm{mod}_{\mathbf{K}}\) 是超度量绝对值，则 \(\mathbf{K}\) 是 \(p\)-进域 \(\mathbf{Q}_p\) 上有限秩的代数。
\item
  若 \(\mathbf{K}\) 的特征为 \(p \neq 0\)，则它同构于一个中心为形式幂级数域 \(\mathbf{F}_q((\mathrm{T}))\) 的域（其中 \(q\) 是 \(p\) 的幂），并且在其中心上有限秩。
\end{enumerate}

\begin{enumerate}
\def\labelenumi{(\roman{enumi})}
\item
由 Ostrowski 定理（§ 6，第 4 节，定理 2），K 是一个同构于 R、C 或 H 的处处稠密子域的拓扑域；又由于 K 完备，它同构于 R、C 或 H。
\item
  令 A 为绝对值 \(mod_{K}\) 的环，m 为其极大理想。我们知道 A/m 是有限域（§ 5，第 1 节，命题 2），因而 \(mod_{K}\) 在 Q 上诱导的绝对值具有有限剩余域；这只有当它等价于 \(p\)-进绝对值时才可能（§ 6，第 3 节，命题 4）。\(\mathbf{Q}\) 在 \(\mathbf{K}\) 中的闭包于是同构于 \(\mathbf{Q}_p\)，并包含于 \(\mathbf{K}\) 的中心中，因为后者在 \(\mathbf{K}\) 中闭；由第 1 节的命题 2 即得结论。
\item
  第二个断言由第 2 节命题 2 的推论和第一个断言得到。为证明第一个断言，注意 \(\mathrm{mod}_{\mathbf{K}}\) 必为超度量绝对值（§ 6，第 2 节，命题 3 的推论）；在第 2 节命题 3 的证明所用的记号中，中心 \(Z\) 是 \(K\) 中与 \(s\) 和 \(u\) 都交换的元素组成的；但是根据 (3)，
\end{enumerate}

\[
u ^ {- q} s u ^ {q} = s ^ {q p ^ {r}} = s,
\]

于是 \(u^q \in \mathbf{Z}\)，我们得出 \(\mathbf{Z}\) 非离散。由于 \(\mathbf{Z}\) 局部紧，我们可限于 \(\mathbf{K}\) 交换的情形。此时的 \(\mathbf{F}_p\)-代数 \(\mathbf{F}_p[s]\) 位于 \(\mathbf{K}\) 中，是有限域，因为 \(s^{q-1} = 1\)，且显然其中每个元素都满足 \(y^q = y\)，所以它与 \(S\) 相同并同构于 \(\mathbf{F}_q\)，因为 \(S \subset \mathbf{F}_p[s]\) 有 \(q\) 个元素。由于 \(S\) 中任意两个元素之和仍在 \(S\) 中，把每个形式幂级数 \(\sum_{i=0}^{\infty} s_i T^i \in \mathbf{F}_q[[T]]\) 映到元素 \(\sum_{i=0}^{\infty}s_{i}u^{i}\) 的映射，是从环 \(F_{q}[[T]]\) 到环 A 的双射同态，从而立即得到结论。

推论 1. 每个非离散局部紧域在其中心上有限秩。

推论 2. 每个局部紧域或者连通，或者全不连通；若它连通，则同构于 R、C 或 H。

若域 K 上的拓扑由超度量绝对值定义，则 K 关于此拓扑全不连通。

注。令 \(s\) 为一个整数 \(>0\)；子域 \(\mathbf{F}_{q}((\mathrm{T}^{s})) = \mathbf{L}\) 是 \(\mathbf{K} = \mathbf{F}_{q}((\mathrm{T}))\) 的子域，在 \(\mathbf{K}\) 中闭，且 \(e(\mathrm{K/L}) = s\)、\(f(\mathrm{K/L}) = 1\)。因此可见，存在闭子域 \(\mathbf{L}\)，它们属于 \(\mathbf{K}\)，非离散且使 \(e(\mathrm{K/L})\)（从而次数 {[}K: L{]}）任意大（这与特征为 0 的局部紧域的情形相反；在那里，每个这样的局部紧子域 \(\mathbf{L}\) 都属于域 \(\mathbf{K}\)，且必包含 \(\mathbf{R}\) 或 \(\mathbf{Q}_{p}\)，因而 {[}K: L{]} 有界）。

